\textbf{Hypotheses and registered tests} (from \texttt{proposal.md}, before the runs). \textbf{H1}: on \texttt{bimodal50} the MSE head has lower success than both GMM and DDPM (Welch/paired tests from \texttt{rh compare}, $\alpha=0.05$). \textbf{H2}: DDPM does not exceed a 2-component GMM by $\ge0.05$ success on \texttt{bimodal50}. \textbf{H3}: on \texttt{bimodal80}, GMM and DDPM reproduce the demonstrated upper-mode share (mean upper\_frac within 0.1 of 0.8). \textbf{H4}: MSE success increases with start jitter, GMM/DDPM stay roughly flat. \textbf{H5}: DDPM success degrades with 1 or 5 denoising steps.

\textbf{Data and protocol.} Each run generates 2,000 demonstrations, trains one head, and evaluates on 200 closed-loop rollouts from fresh start states drawn with a seed-derived generator. Main groups use 10 seeds (0--9: 3 systems $\times$ 3 tasks $=90$ runs); sensitivity and ablation groups use 5 seeds (0--4) on \texttt{bimodal50}. Start jitter is $j=0.01$ in the main runs; this default was chosen from a single pilot seed (99, not reported in tables) in which MSE was neither at 0 nor at 100\% success, so the main task is deliberately in the regime where MSE is partly working; the jitter sweep ($j\in\{0,0.003,0.01,0.03,0.1\}$, the 0.01 point being the main rows) shows the other regimes. In total 216 experiment runs were logged, plus one pilot-seed sanity run; one earlier sanity attempt failed (missing output directory) and was rerun.

\textbf{Hyperparameters and tuning.} No per-method tuning was done, to avoid favouring any head: MLP hidden width 128 (3 hidden layers for MSE and GMM, 4 for DDPM because of its extra inputs, which gives DDPM slightly more capacity), Adam with learning rate $10^{-3}$ and cosine decay, batch 256, 40 epochs, no early stopping and no validation split (nothing is selected). Actions and observations are standardised with training statistics. The DDPM uses 50 steps in the main runs. The pilot seed was used only to debug and to fix the default jitter; no hyperparameter was changed after seeing evaluation results.

\textbf{Compute.} CPU only (2 threads), about 10 s per run; all 216 runs took roughly 25 minutes in total on a shared machine. Results are means $\pm$ standard deviation over seeds; tests are Welch and paired $t$-tests from \texttt{rh compare}, which for 5--10 seeds and saturated metrics (many runs at exactly 1.0) are only indicative.
